% Options for packages loaded elsewhere
\PassOptionsToPackage{unicode}{hyperref}
\PassOptionsToPackage{hyphens}{url}
\documentclass[
]{article}
\usepackage{xcolor}
\usepackage[margin=1in]{geometry}
\usepackage{amsmath,amssymb}
\setcounter{secnumdepth}{-\maxdimen} % remove section numbering
\usepackage{iftex}
\ifPDFTeX
  \usepackage[T1]{fontenc}
  \usepackage[utf8]{inputenc}
  \usepackage{textcomp} % provide euro and other symbols
\else % if luatex or xetex
  \usepackage{unicode-math} % this also loads fontspec
  \defaultfontfeatures{Scale=MatchLowercase}
  \defaultfontfeatures[\rmfamily]{Ligatures=TeX,Scale=1}
\fi
\usepackage{lmodern}
\ifPDFTeX\else
  % xetex/luatex font selection
\fi
% Use upquote if available, for straight quotes in verbatim environments
\IfFileExists{upquote.sty}{\usepackage{upquote}}{}
\IfFileExists{microtype.sty}{% use microtype if available
  \usepackage[]{microtype}
  \UseMicrotypeSet[protrusion]{basicmath} % disable protrusion for tt fonts
}{}
\makeatletter
\@ifundefined{KOMAClassName}{% if non-KOMA class
  \IfFileExists{parskip.sty}{%
    \usepackage{parskip}
  }{% else
    \setlength{\parindent}{0pt}
    \setlength{\parskip}{6pt plus 2pt minus 1pt}}
}{% if KOMA class
  \KOMAoptions{parskip=half}}
\makeatother
\usepackage{color}
\usepackage{fancyvrb}
\newcommand{\VerbBar}{|}
\newcommand{\VERB}{\Verb[commandchars=\\\{\}]}
\DefineVerbatimEnvironment{Highlighting}{Verbatim}{commandchars=\\\{\}}
% Add ',fontsize=\small' for more characters per line
\usepackage{framed}
\definecolor{shadecolor}{RGB}{248,248,248}
\newenvironment{Shaded}{\begin{snugshade}}{\end{snugshade}}
\newcommand{\AlertTok}[1]{\textcolor[rgb]{0.94,0.16,0.16}{#1}}
\newcommand{\AnnotationTok}[1]{\textcolor[rgb]{0.56,0.35,0.01}{\textbf{\textit{#1}}}}
\newcommand{\AttributeTok}[1]{\textcolor[rgb]{0.13,0.29,0.53}{#1}}
\newcommand{\BaseNTok}[1]{\textcolor[rgb]{0.00,0.00,0.81}{#1}}
\newcommand{\BuiltInTok}[1]{#1}
\newcommand{\CharTok}[1]{\textcolor[rgb]{0.31,0.60,0.02}{#1}}
\newcommand{\CommentTok}[1]{\textcolor[rgb]{0.56,0.35,0.01}{\textit{#1}}}
\newcommand{\CommentVarTok}[1]{\textcolor[rgb]{0.56,0.35,0.01}{\textbf{\textit{#1}}}}
\newcommand{\ConstantTok}[1]{\textcolor[rgb]{0.56,0.35,0.01}{#1}}
\newcommand{\ControlFlowTok}[1]{\textcolor[rgb]{0.13,0.29,0.53}{\textbf{#1}}}
\newcommand{\DataTypeTok}[1]{\textcolor[rgb]{0.13,0.29,0.53}{#1}}
\newcommand{\DecValTok}[1]{\textcolor[rgb]{0.00,0.00,0.81}{#1}}
\newcommand{\DocumentationTok}[1]{\textcolor[rgb]{0.56,0.35,0.01}{\textbf{\textit{#1}}}}
\newcommand{\ErrorTok}[1]{\textcolor[rgb]{0.64,0.00,0.00}{\textbf{#1}}}
\newcommand{\ExtensionTok}[1]{#1}
\newcommand{\FloatTok}[1]{\textcolor[rgb]{0.00,0.00,0.81}{#1}}
\newcommand{\FunctionTok}[1]{\textcolor[rgb]{0.13,0.29,0.53}{\textbf{#1}}}
\newcommand{\ImportTok}[1]{#1}
\newcommand{\InformationTok}[1]{\textcolor[rgb]{0.56,0.35,0.01}{\textbf{\textit{#1}}}}
\newcommand{\KeywordTok}[1]{\textcolor[rgb]{0.13,0.29,0.53}{\textbf{#1}}}
\newcommand{\NormalTok}[1]{#1}
\newcommand{\OperatorTok}[1]{\textcolor[rgb]{0.81,0.36,0.00}{\textbf{#1}}}
\newcommand{\OtherTok}[1]{\textcolor[rgb]{0.56,0.35,0.01}{#1}}
\newcommand{\PreprocessorTok}[1]{\textcolor[rgb]{0.56,0.35,0.01}{\textit{#1}}}
\newcommand{\RegionMarkerTok}[1]{#1}
\newcommand{\SpecialCharTok}[1]{\textcolor[rgb]{0.81,0.36,0.00}{\textbf{#1}}}
\newcommand{\SpecialStringTok}[1]{\textcolor[rgb]{0.31,0.60,0.02}{#1}}
\newcommand{\StringTok}[1]{\textcolor[rgb]{0.31,0.60,0.02}{#1}}
\newcommand{\VariableTok}[1]{\textcolor[rgb]{0.00,0.00,0.00}{#1}}
\newcommand{\VerbatimStringTok}[1]{\textcolor[rgb]{0.31,0.60,0.02}{#1}}
\newcommand{\WarningTok}[1]{\textcolor[rgb]{0.56,0.35,0.01}{\textbf{\textit{#1}}}}
\usepackage{graphicx}
\makeatletter
\newsavebox\pandoc@box
\newcommand*\pandocbounded[1]{% scales image to fit in text height/width
  \sbox\pandoc@box{#1}%
  \Gscale@div\@tempa{\textheight}{\dimexpr\ht\pandoc@box+\dp\pandoc@box\relax}%
  \Gscale@div\@tempb{\linewidth}{\wd\pandoc@box}%
  \ifdim\@tempb\p@<\@tempa\p@\let\@tempa\@tempb\fi% select the smaller of both
  \ifdim\@tempa\p@<\p@\scalebox{\@tempa}{\usebox\pandoc@box}%
  \else\usebox{\pandoc@box}%
  \fi%
}
% Set default figure placement to htbp
\def\fps@figure{htbp}
\makeatother
\setlength{\emergencystretch}{3em} % prevent overfull lines
\providecommand{\tightlist}{%
  \setlength{\itemsep}{0pt}\setlength{\parskip}{0pt}}
\usepackage{bookmark}
\IfFileExists{xurl.sty}{\usepackage{xurl}}{} % add URL line breaks if available
\urlstyle{same}
\hypersetup{
  pdftitle={Bonus Question},
  pdfauthor={Hikmat Paudel},
  hidelinks,
  pdfcreator={LaTeX via pandoc}}

\title{Bonus Question}
\author{Hikmat Paudel}
\date{2026-04-09}

\begin{document}
\maketitle

\section{1. Recap Material}\label{recap-material}

Give me notes on the recap material. These questions focus on material
that students commonly struggle with (data wrangling, plotting,
modeling, or debugging).

\#\#Data Wrangling (dplyr)

Data wrangling is the process of preparing and transforming raw data
into a useful format. The \texttt{dplyr} package provides several
important functions:

\begin{itemize}
\tightlist
\item
  \texttt{select()} is used to choose specific columns\\
\item
  \texttt{filter()} extracts rows based on conditions\\
\item
  \texttt{mutate()} creates or modifies variables\\
\item
  \texttt{summarise()} calculates summary statistics\\
\item
  \texttt{group\_by()} groups the data before summarizing
\end{itemize}

It is important to always check column names and ensure variables have
the correct data types before performing operations.

Example

\begin{Shaded}
\begin{Highlighting}[]
\FunctionTok{library}\NormalTok{(dplyr)}
\end{Highlighting}
\end{Shaded}

\begin{verbatim}
## 
## Attaching package: 'dplyr'
\end{verbatim}

\begin{verbatim}
## The following objects are masked from 'package:stats':
## 
##     filter, lag
\end{verbatim}

\begin{verbatim}
## The following objects are masked from 'package:base':
## 
##     intersect, setdiff, setequal, union
\end{verbatim}

\begin{Shaded}
\begin{Highlighting}[]
\CommentTok{\# Create example dataset}
\NormalTok{df }\OtherTok{\textless{}{-}} \FunctionTok{data.frame}\NormalTok{(}
  \AttributeTok{year =} \FunctionTok{c}\NormalTok{(}\DecValTok{2025}\NormalTok{, }\DecValTok{2025}\NormalTok{, }\DecValTok{2024}\NormalTok{, }\DecValTok{2025}\NormalTok{),}
  \AttributeTok{treatment =} \FunctionTok{c}\NormalTok{(}\StringTok{"A"}\NormalTok{, }\StringTok{"B"}\NormalTok{, }\StringTok{"A"}\NormalTok{, }\StringTok{"B"}\NormalTok{),}
  \AttributeTok{yield =} \FunctionTok{c}\NormalTok{(}\DecValTok{10}\NormalTok{, }\DecValTok{15}\NormalTok{, }\DecValTok{8}\NormalTok{, }\DecValTok{20}\NormalTok{),}
  \AttributeTok{rep1 =} \FunctionTok{c}\NormalTok{(}\DecValTok{5}\NormalTok{,}\DecValTok{6}\NormalTok{,}\DecValTok{4}\NormalTok{,}\DecValTok{7}\NormalTok{),}
  \AttributeTok{rep2 =} \FunctionTok{c}\NormalTok{(}\DecValTok{7}\NormalTok{,}\DecValTok{5}\NormalTok{,}\DecValTok{6}\NormalTok{,}\DecValTok{8}\NormalTok{)}
\NormalTok{)}

\CommentTok{\# Perform filtering, grouping, and summarizing}
\NormalTok{data\_summary }\OtherTok{\textless{}{-}}\NormalTok{ df }\SpecialCharTok{\%\textgreater{}\%}
  \FunctionTok{filter}\NormalTok{(year }\SpecialCharTok{==} \DecValTok{2025}\NormalTok{) }\SpecialCharTok{\%\textgreater{}\%}
  \FunctionTok{group\_by}\NormalTok{(treatment) }\SpecialCharTok{\%\textgreater{}\%}
  \FunctionTok{summarise}\NormalTok{(}\AttributeTok{mean\_yield =} \FunctionTok{mean}\NormalTok{(yield, }\AttributeTok{na.rm =} \ConstantTok{TRUE}\NormalTok{))}

\CommentTok{\# Show results}
\NormalTok{data\_summary}
\end{Highlighting}
\end{Shaded}

\begin{verbatim}
## # A tibble: 2 x 2
##   treatment mean_yield
##   <chr>          <dbl>
## 1 A               10  
## 2 B               17.5
\end{verbatim}

Common Mistakes Using summarise() without grouping first Incorrect or
misspelled column names Treating categorical variables as numeric

\#\#Plotting (ggplot2) The ggplot2 package is used to create
visualizations. A basic plot structure looks like this:

\begin{Shaded}
\begin{Highlighting}[]
\FunctionTok{library}\NormalTok{(ggplot2)}
\end{Highlighting}
\end{Shaded}

\begin{verbatim}
## Warning: package 'ggplot2' was built under R version 4.5.2
\end{verbatim}

\begin{Shaded}
\begin{Highlighting}[]
\CommentTok{\# Scatter plot with real columns}
\FunctionTok{ggplot}\NormalTok{(df, }\FunctionTok{aes}\NormalTok{(}\AttributeTok{x =}\NormalTok{ year, }\AttributeTok{y =}\NormalTok{ yield)) }\SpecialCharTok{+}
  \FunctionTok{geom\_point}\NormalTok{() }\SpecialCharTok{+}
  \FunctionTok{theme\_classic}\NormalTok{()}
\end{Highlighting}
\end{Shaded}

\pandocbounded{\includegraphics[keepaspectratio]{Bonus-material_files/figure-latex/unnamed-chunk-2-1.pdf}}
Common layers include: geom\_point() for scatterplots geom\_line() for
line graphs geom\_bar(stat = ``identity'') for bar charts facet\_wrap()
for multiple panels theme\_classic() for a simple design

Example

\begin{Shaded}
\begin{Highlighting}[]
\FunctionTok{ggplot}\NormalTok{(df, }\FunctionTok{aes}\NormalTok{(}\AttributeTok{x =}\NormalTok{ treatment, }\AttributeTok{y =}\NormalTok{ yield)) }\SpecialCharTok{+}
  \FunctionTok{geom\_boxplot}\NormalTok{() }\SpecialCharTok{+}
  \FunctionTok{theme\_classic}\NormalTok{()}
\end{Highlighting}
\end{Shaded}

\pandocbounded{\includegraphics[keepaspectratio]{Bonus-material_files/figure-latex/unnamed-chunk-3-1.pdf}}

Common Mistakes Incorrect placement of variables inside/outside aes()
Forgetting to use + between layers Using wrong variable names Plotting
unsummarized data incorrectly

\#\#Linear Modeling Linear models are used to study relationships
between variables.

Examples

\begin{Shaded}
\begin{Highlighting}[]
\NormalTok{model }\OtherTok{\textless{}{-}} \FunctionTok{lm}\NormalTok{(yield }\SpecialCharTok{\textasciitilde{}}\NormalTok{ treatment, }\AttributeTok{data =}\NormalTok{ df)}
\FunctionTok{summary}\NormalTok{(model)}
\end{Highlighting}
\end{Shaded}

\begin{verbatim}
## 
## Call:
## lm(formula = yield ~ treatment, data = df)
## 
## Residuals:
##    1    2    3    4 
##  1.0 -2.5 -1.0  2.5 
## 
## Coefficients:
##             Estimate Std. Error t value Pr(>|t|)  
## (Intercept)    9.000      1.904   4.727   0.0420 *
## treatmentB     8.500      2.693   3.157   0.0874 .
## ---
## Signif. codes:  0 '***' 0.001 '**' 0.01 '*' 0.05 '.' 0.1 ' ' 1
## 
## Residual standard error: 2.693 on 2 degrees of freedom
## Multiple R-squared:  0.8329, Adjusted R-squared:  0.7493 
## F-statistic: 9.966 on 1 and 2 DF,  p-value: 0.08739
\end{verbatim}

Important Outputs Estimate: shows the effect size p-value: indicates
statistical significance R-squared: measures model fit Residuals: show
prediction errors

Prediction

\begin{Shaded}
\begin{Highlighting}[]
\FunctionTok{predict}\NormalTok{(model)}
\end{Highlighting}
\end{Shaded}

\begin{verbatim}
##    1    2    3    4 
##  9.0 17.5  9.0 17.5
\end{verbatim}

Common Mistakes Mixing up predictor and response variables Not
converting categorical variables into factors Ignoring missing values in
the dataset

\#\#Debugging Debugging helps identify and fix errors in code. Useful
strategies include: Carefully reading error messages Running code
step-by-step Checking object and variable names Verifying parentheses
and commas Inspecting data using str() and names()

Example

\begin{Shaded}
\begin{Highlighting}[]
\FunctionTok{str}\NormalTok{(df)}
\end{Highlighting}
\end{Shaded}

\begin{verbatim}
## 'data.frame':    4 obs. of  5 variables:
##  $ year     : num  2025 2025 2024 2025
##  $ treatment: chr  "A" "B" "A" "B"
##  $ yield    : num  10 15 8 20
##  $ rep1     : num  5 6 4 7
##  $ rep2     : num  7 5 6 8
\end{verbatim}

\begin{Shaded}
\begin{Highlighting}[]
\FunctionTok{names}\NormalTok{(df)}
\end{Highlighting}
\end{Shaded}

\begin{verbatim}
## [1] "year"      "treatment" "yield"     "rep1"      "rep2"
\end{verbatim}

\#2. Extra Topic Exploration

Complete one of the extra optional materials provided, OR explore a
topic of your choice, or find a topic online and explore the code and
how it works.

I explored the pivot\_longer() function from the tidyr package, which is
used to reshape data from wide format into long format. This is
especially useful when preparing data for visualization and analysis.

Example

\begin{Shaded}
\begin{Highlighting}[]
\FunctionTok{library}\NormalTok{(tidyr)}
\end{Highlighting}
\end{Shaded}

\begin{verbatim}
## Warning: package 'tidyr' was built under R version 4.5.2
\end{verbatim}

\begin{Shaded}
\begin{Highlighting}[]
\NormalTok{long\_df }\OtherTok{\textless{}{-}}\NormalTok{ df }\SpecialCharTok{\%\textgreater{}\%}
  \FunctionTok{pivot\_longer}\NormalTok{(}
    \AttributeTok{cols =} \FunctionTok{starts\_with}\NormalTok{(}\StringTok{"rep"}\NormalTok{),}
    \AttributeTok{names\_to =} \StringTok{"replicate"}\NormalTok{,}
    \AttributeTok{values\_to =} \StringTok{"value"}\NormalTok{)}

\NormalTok{long\_df}
\end{Highlighting}
\end{Shaded}

\begin{verbatim}
## # A tibble: 8 x 5
##    year treatment yield replicate value
##   <dbl> <chr>     <dbl> <chr>     <dbl>
## 1  2025 A            10 rep1          5
## 2  2025 A            10 rep2          7
## 3  2025 B            15 rep1          6
## 4  2025 B            15 rep2          5
## 5  2024 A             8 rep1          4
## 6  2024 A             8 rep2          6
## 7  2025 B            20 rep1          7
## 8  2025 B            20 rep2          8
\end{verbatim}

Why It Is Useful Makes plotting easier Simplifies grouping and
summarizing Helps maintain a tidy data structure

Justification I chose this topic because reshaping data is an important
step in many analyses. I wanted to better understand how to organize
datasets in a way that works efficiently with visualization and modeling
tools. This will help me work with more complex datasets in the future.

\section{3. Introduction to renv}\label{introduction-to-renv}

Follow the coding demonstration in class. Turn in notes.

\subsection{renv Notes}\label{renv-notes}

\textbf{Purpose:} - Creates reproducible R environments\\
- Keeps track of package versions\\
- Helps others run the same code without issues

\subsection{Initialize Project}\label{initialize-project}

\begin{Shaded}
\begin{Highlighting}[]
\FunctionTok{install.packages}\NormalTok{(}\StringTok{"renv"}\NormalTok{)}
\end{Highlighting}
\end{Shaded}

\begin{verbatim}
## The following package(s) will be installed:
## - renv [1.2.0]
## These packages will be installed into "C:/Users/admin/Desktop/Assignments-repository/renv/library/windows/R-4.5/x86_64-w64-mingw32".
## 
## # Installing packages --------------------------------------------------------
## ✔ renv 1.2.0                               [linked from cache]
## Successfully installed 1 package in 0.21 seconds.
\end{verbatim}

\begin{Shaded}
\begin{Highlighting}[]
\NormalTok{renv}\SpecialCharTok{::}\FunctionTok{init}\NormalTok{()}
\end{Highlighting}
\end{Shaded}

\begin{verbatim}
## - The project is out-of-sync -- use `]8;;x-r-run:renv::status()renv::status()]8;;` for details.
\end{verbatim}

Install Packages

\begin{Shaded}
\begin{Highlighting}[]
\FunctionTok{install.packages}\NormalTok{(}\StringTok{"dplyr"}\NormalTok{)}
\end{Highlighting}
\end{Shaded}

\begin{verbatim}
## The following package(s) will be installed:
## - dplyr [1.2.1]
## These packages will be installed into "C:/Users/admin/Desktop/Assignments-repository/renv/library/windows/R-4.5/x86_64-w64-mingw32".
## 
## # Installing packages --------------------------------------------------------
## ✔ dplyr 1.2.1                              [linked from cache]
## Successfully installed 1 package in 16 milliseconds.
\end{verbatim}

Save Environment

\begin{Shaded}
\begin{Highlighting}[]
\NormalTok{renv}\SpecialCharTok{::}\FunctionTok{snapshot}\NormalTok{()}
\end{Highlighting}
\end{Shaded}

\begin{verbatim}
## The following required packages are not installed:
## - kable
## Packages must first be installed before renv can snapshot them.
## Use `]8;;x-r-run:renv::dependencies()renv::dependencies()]8;;` to see where this package is used in your project.
## 
## - The lockfile is already up to date.
\end{verbatim}

Restore Environment

\begin{Shaded}
\begin{Highlighting}[]
\NormalTok{renv}\SpecialCharTok{::}\FunctionTok{restore}\NormalTok{()}
\end{Highlighting}
\end{Shaded}

\begin{verbatim}
## - The library is already synchronized with the lockfile.
\end{verbatim}

Files Created renv.lock → stores package version details renv/ folder →
contains the project library

Why It Is Important Ensures results can be reproduced Avoids conflicts
between package versions Makes it easier to share projects with others

\end{document}
